Lãi suất phi rủi ro $r_f$ là tỷ suất sinh lời của một tài sản được coi là không có rủi ro vỡ nợ, thường là trái phiếu chính phủ ngắn hạn. Trong các mô hình định giá, lãi suất này được dùng làm lợi nhuận cơ sở để tính toán phần bù rủi ro. Khi kết hợp tài sản phi rủi ro với một danh mục tài sản rủi ro, tập hợp các danh mục đầu tư khả thi sẽ tạo thành một đường thẳng nối từ điểm phi rủi ro và đi qua điểm danh mục tài sản rủi ro đó. Đường biên hiệu quả mới lúc này không còn là hình hyperbol mà trở thành một đường thẳng tiếp tuyến với đường biên hiệu quả của các tài sản rủi ro. Đường này được gọi là đường thị trường vốn (Capital Market Line - CML).